% !PRINTED 64
\weberpage{64}
%<<CONT
$\nu$ durch $\nu + \nu'$ ersetzt, wenn $\nu'$ den Index von $n'$ bedeutet.
Dadurch folgt aus (3):
\[
\varepsilon^{-\beta \nu'} (\varepsilon^{\beta}, r) = \sum^{n} \varepsilon^{\beta \nu} r^{n n'},
\]
und wenn man also mit $r^{n'}$ multiplicirt und die Summe \"uber $n'$
nimmt:
\weq{4}{(\varepsilon^{-\beta}, r)(\varepsilon^{\beta}, r) = \sum^{n, n'} \varepsilon^{\beta \nu} r^{n'(n+1)}.}
Es ist also zun\"achst die Summe
\weq{5}{\sigma = \sum^{n'} r^{n'(n+1)}}
zu bestimmen. Die Summe $\sigma$ verschwindet aber immer, wenn
$n + 1$ nicht durch $q^{\varkappa - 1}$ theilbar ist (nach Bd.\ I, \S.\,133); denn
dann ist $r^{n+1}$ eine Einheitswurzel, deren Grad eine h\"ohere als
die erste Potenz von $q$ ist. $\sigma$ kann also nur dann von Null verschieden
 sein, wenn $n$ die Form hat
\weq{6}{n \equiv -1 + t q^{\varkappa - 1} \pmod{m},}
und dann wollen wir seinen Werth mit $\sigma_t$ bezeichnen. Wir erhalten
 alle Werthe von $n$ nach dem Modul $m$, die in der Form
(6) enthalten sind, wenn wir $t$ ein volles Restsystem nach dem
Modul $q$ durchlaufen lassen, also etwa
\weq{7}{t = 0, 1, 2 \ldots q - 1}
setzen. Die Summe $\sigma$ besteht aus $\varphi(m)$ Gliedern. Ist $t = 0$, so
wird jedes dieser Glieder $= 1$ und es folgt
\weq{8}{\sigma_0 = \varphi(m) = q^{\varkappa} - q^{\varkappa - 1};}
f\"ur jeden anderen Werth von $t$ ist $r^{n+1}$ eine primitive $q^{\text{te}}$ Einheitswurzel
, und in $\sigma_t$ kommt dann jede solche Einheitswurzel
$\varphi(m) : (q - 1) = q^{\varkappa - 1}$ mal vor. Die Summe aller primitiven
$q^{\text{ten}}$ Einheitswurzeln ist aber gleich $-1$ und also folgt f\"ur
$t = 1, 2 \ldots q - 1$
\weq{9}{\sigma_t = -q^{\varkappa - 1}.}
Damit ist die Summe $\sigma$ bestimmt.

Um aber den Werth der Summe in (4) daraus abzuleiten,
ist es noch n\"othig, den Index $\nu$ der Zahlen $n$ von der Form (6)
zu ermitteln.

F\"ur diese Zahlen ist aber
\[
g^{\nu} \equiv -1 \pmod{q^{\varkappa - 1}},
\]
und da nach dem \gesperrt{Fermat'schen} Satze [\S.\,14, (11)]
\weq{10}{g^{1/2 \varphi(q^{\varkappa})} \equiv -1 \pmod{q^{\varkappa}}}
%CONT>>
